\chapter{Introduction to the Project}

\section{Overview}

In recent years, advancements in artificial intelligence and deep learning have significantly improved the performance of computer vision systems. Image classification, a cornerstone of these systems, has seen remarkable progress with the advent of Convolutional Neural Networks (CNNs). However, conventional deep learning models predominantly rely on spatial domain features—the raw pixel values—often overlooking the rich, discriminative patterns embedded within the frequency domain. This limitation can restrict model robustness and interpretability, particularly in complex visual tasks such as generalized object recognition and precise facial analysis.

While standard CNN architectures like ResNet \cite{resnet} have demonstrated strong performance on spatial patterns, they can struggle to explicitly capture periodic textures and structural geometries that are highly visible in the spectral domain. Recent studies in frequency learning for image classification \cite{freq_learning} highlight the need for systems that explicitly leverage these spectral representations to ensure both robustness and improved feature discrimination.

This project proposes a novel, hybrid image classification pipeline that explicitly leverages both spatial and frequency domain representations. By integrating Fast Fourier Transform (FFT) techniques \cite{fft} with state-of-the-art machine learning algorithms, the proposed framework extracts both global and blockwise spectral features. These features provide a highly comprehensive understanding of image textures prior to classification.

The proposed framework adopts a multi-component architecture consisting of advanced data preprocessing (Histogram Equalization and Gaussian Blur), frequency feature extraction, Principal Component Analysis (PCA) \cite{pca} for dimensionality reduction, and robust classifiers including XGBoost \cite{xgboost} and PyTorch-based neural networks \cite{pytorch}. Furthermore, the integration of visual explanation mechanisms like Grad-CAM \cite{gradcam} ensures that the model's decision-making is localized to semantically meaningful regions, thereby improving interpretability.

\section{Problem Statement}

The widespread adoption of deep learning in image classification has led to models that achieve high accuracy but often act as "black boxes" that rely heavily on local spatial textures. Existing conventional approaches operate primarily on raw pixels and can fail to generalize well when spatial textures are ambiguous or corrupted. Therefore, there is a need for a classification system that can inherently understand the global and local structural frequencies of an image to improve both accuracy and reliability.

The key challenge is to design an architecture that effectively extracts and fuses frequency domain features with standard spatial representations without introducing overwhelming computational overhead. These requirements are difficult to achieve using standard, out-of-the-box CNN architectures which are heavily biased towards spatial pooling.

This project addresses this challenge by engineering complementary feature sets—one relying on standard deep spatial representations (capable of high-level semantic understanding) and another relying on explicit frequency domain transformations (highly sensitive to periodic textures and underlying structure).

Existing deep learning approaches often require massive datasets to implicitly learn frequency behaviors. In contrast, explicitly supplying these features using analytical methods like the Fast Fourier Transform provides a more direct and efficient learning path. These limitations in current CNN paradigms motivate the need for a hybrid framework that balances spatial and spectral learning for effective image classification.

\section{Objectives}

The primary objective of this project is to design a hybrid, frequency-aware image classification framework that ensures high accuracy, robustness, and interpretability in computer vision tasks. The system aims to extract explicitly engineered frequency features—both global and blockwise—that complement traditional deep learning representations. 

Specific objectives include:
\begin{itemize}
    \item Integrating Fast Fourier Transform (FFT) techniques for robust feature extraction across both global images and localized blocks.
    \item Implementing dimensionality reduction via Principal Component Analysis (PCA) to streamline complex spectral data for classical machine learning models.
    \item Training and evaluating robust classifiers, including XGBoost and custom PyTorch neural networks (such as FFTNet and ResNet18 baselines).
    \item Validating the interpretability of the models using Grad-CAM to ensure decisions are based on semantically meaningful features rather than spurious artifacts.
\end{itemize}

\section{Scope of the Project}

This project covers the development of a proactive, hybrid image classification system, including dataset preparation, modular architecture design, training pipelines, and comprehensive evaluation. The system operates on standard benchmark datasets, specifically STL-10 for generalized object recognition and VGGFace for complex facial analysis. 

The technical scope includes the implementation of preprocessing steps, feature extraction engines, and model training scripts utilizing PyTorch, XGBoost, and scikit-learn. The pipeline is designed for scalability and remote multi-GPU execution. The scope also includes evaluating the framework using standard classification metrics (such as accuracy and cross-entropy loss) and visual explainability metrics. Real-time edge deployment and large-scale web server implementation are beyond the scope of this work.
